\begin{problem}[第35届全国中学生物理竞赛决赛][准粒子动力学]{102}
    在固体材料中，考虑相互作用后，可以利用“准粒子”的概念研究材料的物理性质。准粒子的动能与动量之间的关系可能与真实粒子的不同。当外加电场或磁场时，准粒子的运动往往可以用经典力学的方法来处理。在某种二维界面结构中，存在电量为 $q$ 、有效质量为 $m$ 的准粒子，它只能在 $x - y$ 平面内运动，其动能 $K$ 与动量大小 $p$ 之间的关系可表示为 $K = \frac{p^2}{2m} + \alpha p$ ，其中 $\alpha$ 为正的常量。
    \begin{subproblem}
        对于质量为 $m$ 的真实的自由粒子，动能 $K$ 与动量大小 $p$ 之间的关系可表示为 $K = \frac{p^{2}}{2m}$ ，试从动能定理出发，推导该粒子运动的速度 $v$ 与动量 $p$ 之间的关系式；
    \end{subproblem}
    \begin{subproblem}
        仿照（1）的方法，推导准粒子运动的速度 $v$ 与动量 $p$ 之间的关系式；
    \end{subproblem}
    \begin{subproblem}
        用动能表示准粒子运动速度的大小;
    \end{subproblem}
    \begin{subproblem}
        将该二维界面结构置于匀强磁场中，磁场沿 $z$ 轴正方向，磁感应强度大小为 $B$ ，求动能为 $K$ 的准粒子做匀速率圆周运动的半径、周期和角动量的大小；
    \end{subproblem}
    \begin{subproblem}
        将该二维界面结构放置在匀强电场中，准粒子可能在垂直于电场的方向上产生加速度。如果电场沿 $x$ 轴正方向，电场强度大小为 $E$ 。当准粒子的速度大小为 $v$ （ $v \neq \alpha$ ）、方向与 $x$ 轴正方向成 $\theta$ 角时，求其运动的加速度的分量 $a_x$ 和 $a_y$ 。
    \end{subproblem}
\end{problem}